\chapter{Conclusion}

\label{ch:conclusion}

\section{Summary}

\label{sec:conclusion-summary}

This thesis examined where a frozen-model text-to-SPARQL pipeline over Wikidata fails and whether targeted interventions can reduce those failures.
It uses \method{}, which keeps every stage fixed except the one under investigation and evaluates that stage on its own.
\Cref{tab:conclusion-summary} collects the headline figures for each research question.
The paragraphs below state the findings, and the results and discussion chapters give the full evidence and its limits.

\begin{table}[htbp]
\centering
\small
\caption{Findings by research question. Evidence, intervals, caveats are in Results and Discussion.}
\label{tab:conclusion-summary}
\begin{tabular}{@{}p{0.9cm}p{7.6cm}p{4.4cm}@{}}
\toprule
RQ & Finding & Key figures \\
\midrule
1 & Frozen models parse and execute reliably but are rarely correct; failures are shared across models & 24.5 best frozen baseline; 307/400 solved by none of four models \\
2 & Under gold labels and near-ceiling candidate coverage, disambiguation is not the binding constraint; query generation is & 97.4 gold-label vs.\ 30.9 end-to-end (66.5-pt gap) \\
3 & Reasoning over retrieved candidates outperforms specialised linkers; most of its accuracy comes from findable labels & +24.9 [22.3, 27.7] pts over GLiNKER on the same candidates; +4.1 [2.7, 5.6] over the first search result \\
4 & General advice gives small gains; graph knowledge helps as a restriction, not as evidence & $-$14.65 pts (evidence) vs.\ $+$18.99 pts (closed vocabulary) \\
5 & Some failures are detectable from execution alone, without gold supervision & +4.83 [2.93, 7.08] pts; no row can worsen by design \\
\bottomrule
\end{tabular}
\end{table}

\paragraph{RQ1.}

Frozen frontier models write SPARQL that executes but is rarely correct, and their failures are shared across models.

\paragraph{RQ2.}

With gold-derived labels and near-ceiling candidate coverage, the binding constraint is query generation, not disambiguation.
Supplying the concepts helps more than naming the constructs, but composing the query around them remains the largest step: with both supplied, accuracy is 50.4 against 97.4 when the composed reference query is given.

\paragraph{RQ3.}

Reasoning over retrieved candidates clearly outperforms specialised end-to-end linkers, and evaluation asymmetry alone does not explain the gap.
Disambiguation hardly depends on the backbone, while label generation does.

\paragraph{RQ4.}

General advice gives small gains, and graph knowledge helps when it restricts the generator's choices.

\paragraph{RQ5.}

Some of the remaining failures can be detected from execution alone, without gold supervision.
The execution-guided cascade recovers about two fifths of the available oracle headroom.

Together, these findings yield a working pipeline in which every step is chosen on experimental evidence, and a protocol for finding where such pipelines fail.

\section{The Resulting System}

\label{sec:conclusion-system}

The measurements settle several design choices that are often made by convention, and together they describe a configuration that others can start from.

Label generation needed the strongest model, while for disambiguation an open 7B model came within two points of it (95.8 against 97.7 entity F1).
Scale matters at only one stage: the 14B open model produces unfindable labels at 4.9 times the frontier rate, whereas disambiguation varies by less than three percentage points across model sizes.

For linking, retrieving candidates and reasoning over them worked better than linking end to end.
Both fine-tuned linkers fail mainly because they detect mentions in the question text and therefore miss identifiers that the question never names (\cref{sec:res-linkers}).
On a second benchmark, reasoning over candidates again links entities accurately (90.6\,\% entity accuracy), so the result is not specific to one dataset.

Constrain the vocabulary; do not enumerate the schema.
In this pipeline, graph knowledge supplied before generation helped only when it removed options.

An empty result is a useful signal at inference time, because structurally wrong queries often return nothing.
The cascade built on this signal requires no training or gold answers, only replaces failed queries, and should be the baseline for any learned selector.

Candidate pools need to be checked before any linking number can be trusted.
Fixing silently truncated pools raised the gold-label condition by 26 points, because a retrieval layer that returns an empty list because of throttling looks exactly like one that returns an empty list because nothing matches.

The open problem is building a query around a construct that the model can identify but does not yet use correctly.
The next section describes the research directions that follow from this.

\section{Future Work}

\label{sec:future-work}

\paragraph{Grounding by constraint through
filtering.}

The prompt experiment here tests the framing side of this principle at the prompt level, whereas \citet{zhang2026saga} remove incompatible candidates mechanically.
Such filtering replaces the model's obedience with the reliability of Wikidata's own assertions.
The constraint-filtering experiment shows, however, that these assertions can be weak: filtering on every declared type constraint removes the reference property for 7.5\,\% of the mentions it affects, while filtering only on mandatory constraints is safe but does almost nothing.
Whether the safe filter improves accuracy is still open.

\paragraph{Closing the open model's vocabulary
gap.}

String normalisation is not enough: it repairs 43 of 270 unfindable labels, because the main problem is often a different property name rather than a malformed one.
Supplying the vocabulary through constrained generation or light adaptation could make the hybrid deployment discussed earlier unnecessary.

\paragraph{A learned selector over diverse
generations.}

About seven points of oracle headroom remain beyond the cascade.
A learned selector should therefore be evaluated against the cascade, not only against the best single run.

\paragraph{Extending the ground-truth
derivation.}

The linking ground-truth derivation used here needs only aligned gold and label-form queries.
LC-QuAD 2.0~\parencite{dubey2019lcquad2} was used here only to replicate the linking measurement; a full stage-wise analysis on it, including the structural gap, would show whether the attribution holds beyond Instruct-to-SPARQL.
The Wikidata Query Logs dataset~\parencite{walter2026wdql} would test whether it holds for questions distributed like real query logs.

\paragraph{Assessing whether reference noise favours imitation over reasoning}

The discussion identified a mechanism worth testing directly: a fine-tuned system reproduces reference queries (56\,\% character-identical), so an unsound reference rewards imitation over reasoning, and part of the measured gap may come from the answer key itself.
Testing this requires extending the reference-query adjudication to the official split, since nine overlapping rows are not enough.

\paragraph{Re-adjudicating the benchmark's reference
queries.}

Given the 24.6\,\% reference-noise estimate, a large-scale correction should judge the query text rather than the executed output, use two readers with reconciliation, and also sample queries that a heuristic did not flag, since every correction in this study came from that group.
